\clearpage
\pdfbookmark[2]{3. Determination of the Elements}{parabolisch-003}
\begin{center}
[3.]\par
\medskip
On parabolic motion.\par
\medskip
[Determination of the elements from two radius vectors \(r\), \(r'\) and the chord \(\rho\).]
\end{center}

Let
\[
\mathrm{[1]}\qquad \frac{r'-r}{\rho}=\sin\psi,\qquad \frac{\rho}{r'+r}=\sin\varphi.
\]
\SourceLeaf{334}[Using the notation of the Theoria motus, Articles 98 and 108, we have
\[
\begin{gathered}
\frac{p(r'+r)}{2rr'}=1+\cos f\cos F,\qquad
\frac{p(r'-r)}{2rr'}=\sin f\sin F\\
\frac{pp}{rr'}=(\cos f+\cos F)^2;
\end{gathered}
\]
hence
\[
\begin{gathered}
\frac{r'+r}{2p}=\frac{1+\cos f\cos F}{(\cos f+\cos F)^2}\mathpunct{\raisebox{.45ex}{,}}\\
\frac{r'-r}{2p}=\frac{\sin f\sin F}{(\cos f+\cos F)^2}
\end{gathered}
\]
or, since
\[
p=\frac{\rho}{\eta(\theta'-\theta)}\mathpunct{\raisebox{.45ex}{,}}
\]
where
\[
\begin{gathered}
\eta\eta=1+\tfrac14(\theta'+\theta)^2=\frac{1+\cos f^2+2\cos f\cos F}{(\cos f+\cos F)^2}\mathpunct{\raisebox{.45ex}{,}}\\
\theta'-\theta=\frac{2\sin f}{\cos f+\cos F}\mathpunct{\raisebox{.45ex}{,}}
\end{gathered}
\]
we obtain
\[
\begin{gathered}
\frac{\rho}{r'+r}=\sin\varphi=\frac{\sin f\BareRoot{(1+\cos f^2+2\cos f\cos F)}}{1+\cos f\cos F}\\
\frac{r'-r}{\rho}=\sin\psi=\frac{\sin F}{\BareRoot{(1+\cos f^2+2\cos f\cos F)}}\mathpunct{\cdot}
\end{gathered}
\]

From this we obtain
\[
\begin{gathered}
\cos\varphi=\frac{\cos f(\cos f+\cos F)}{1+\cos f\cos F}\mathpunct{\raisebox{.45ex}{,}}\\
\cos\psi=\frac{\cos f+\cos F}{\BareRoot{(1+\cos f^2+2\cos f\cos F)}}\mathpunct{\raisebox{.45ex}{,}}
\end{gathered}
\]
hence]
\[
\mathrm{[2]}\qquad \cos\psi\,\mathrm{tang}\,\varphi=\mathrm{tang}\,f
\]
[from which \(f\) is found.

Furthermore,
\[
\begin{gathered}
\sin\varphi\sin\psi=\frac{\sin f\sin F}{1+\cos f\cos F}\\
\sin\varphi\cos\psi=\frac{\sin f(\cos f+\cos F)}{1+\cos f\cos F}\mathpunct{\raisebox{.25ex}{;}}
\end{gathered}
\]
from which we obtain
\[
\begin{gathered}
\sin\varphi\sin(F-\psi)=\cos f\sin\varphi\sin\psi=\sin f\,\mathrm{tang}\,\psi\cos\varphi\\
\sin\varphi\cos(F-\psi)=\sin f.
\end{gathered}
\]

These equations themselves, or the following equations derived from them,]
\[
\mathrm{tang}\,\psi\cos\varphi=\mathrm{tang}\,(F-\psi)
\]
or
\[
\begin{gathered}
\sin f\sin\psi\,\mathrm{tang}\,\tfrac12\varphi=\sin(2\psi-F)\\
\sin f\sin\psi\,\mathrm{cotang}\,\tfrac12\varphi=\sin F
\end{gathered}
\]
\SourceLeaf{335}[serve to determine \(F\).] The true anomalies [are then] \(F-f\) and \(F+f\).

[We also have
\[
\frac{r'+r}{2p}=\frac{\sin f}{\sin F}\frac{\mathrm{tang}\,\psi}{\sin\varphi\cos\psi}=\frac{\mathrm{tang}\,\tfrac12\varphi}{\sin\varphi\cos\psi^2}\mathpunct{\raisebox{.45ex}{,}}
\]
hence]
\[
q=\left[\frac p2=\right]\tfrac12(r'+r)\cos\tfrac12\varphi^2\,.\cos\psi^2,
\]
[from which \(q\) is known.

To calculate the times elapsed since the passage through perihelion, we have, in the notation of Note 1,
\[
\begin{aligned}
mm&=2(r'+r)\cos(45^0-\tfrac12\varphi)^2\\
nn&=2(r'+r)\sin(45^0-\tfrac12\varphi)^2\\
mn&=(r'+r)\cos\varphi
\end{aligned}
\]
hence
\[
\begin{aligned}
m+n&=2\BareRoot{(r'+r)}\,.\cos\tfrac12\varphi\\
m-n&=2\BareRoot{(r'+r)}\,.\sin\tfrac12\varphi.
\end{aligned}
\]
Now, since
\[
r'-r=(r'+r)\sin\psi\sin\varphi,
\]
we obtain]
\[
\tfrac12(t'+t)=\frac{(r'+r)^{\frac32}\sin\psi\cos\tfrac12\varphi}{6k}(3-2\sin\psi^2\cos\tfrac12\varphi^2).
\]
\begin{center}\rule{.3\linewidth}{.4pt}\end{center}
